\documentclass[11pt,a4paper]{article}
\usepackage[T1]{fontenc}
\usepackage[utf8]{inputenc}
\usepackage{lmodern,amsmath,amssymb,amsthm,mathtools}
\usepackage[margin=27mm]{geometry}
\usepackage{microtype,enumitem,booktabs,longtable,array}
\usepackage{xcolor}
\usepackage[colorlinks=true,linkcolor=blue!45!black,citecolor=blue!45!black,urlcolor=blue!45!black]{hyperref}
\usepackage{fancyhdr}
\pagestyle{fancy}\fancyhf{}
\fancyhead[L]{\small Delayed NNLIF periodic branch}
\fancyhead[R]{\small Matthew J. Colbrook}
\fancyfoot[C]{\thepage}
\setlength{\headheight}{14pt}
\setlength{\parskip}{4pt}
\setlength{\parindent}{0pt}
\numberwithin{equation}{section}
\newtheorem{theorem}{Theorem}[section]
\newtheorem{lemma}[theorem]{Lemma}
\newtheorem{proposition}[theorem]{Proposition}
\newtheorem{corollary}[theorem]{Corollary}
\theoremstyle{remark}\newtheorem{remark}[theorem]{Remark}
\newcommand{\R}{\mathbb R}\newcommand{\C}{\mathbb C}\newcommand{\T}{\mathbb T}
\newcommand{\HH}{\mathcal H}\newcommand{\DD}{\mathcal D}
\newcommand{\dd}{\,\mathrm d}
\DeclareMathOperator{\supp}{supp}
\DeclareMathOperator{\Rea}{Re}
\DeclareMathOperator{\Ima}{Im}
\DeclareMathOperator{\Ker}{ker}
\DeclareMathOperator{\Ran}{ran}
\DeclareMathOperator{\PV}{PV}
\emergencystretch=2em
\hypersetup{pdftitle={A periodic branch for the delayed noisy leaky integrate-and-fire equation},pdfauthor={Matthew J. Colbrook}}
\title{\textbf{A periodic branch for the delayed\\noisy leaky integrate-and-fire equation}}
\author{Matthew J. Colbrook\\[3pt]
\small Department of Applied Mathematics and Theoretical Physics\\
\small University of Cambridge, Cambridge CB3 0WA, United Kingdom\\
\small \href{mailto:m.colbrook@damtp.cam.ac.uk}{m.colbrook@damtp.cam.ac.uk}}
\date{2 October 2026}
\begin{document}
\maketitle
\begin{abstract}
We propose a local branch of nonnegative periodic solutions of the delayed noisy leaky integrate-and-fire reset equation with inhibitory coupling and nonconstant firing rate. Translating both voltage thresholds with the coupling keeps a centered stationary reset profile fixed. An explicit boundary-response function gives a nondegenerate two-parameter reduction at high frequency. Resolvent bounds, time-Sobolev trace estimates, a negative-part energy argument and parameter reconstruction address the full PDE on its voltage half-line. The result uses the repository target's freedom to choose negative physical thresholds.
\end{abstract}
\smallskip
\noindent\textbf{Submission.} AIM 353 at repository revision \texttt{fa98b7525fa3}; solution claim awaiting pull-request review.

\section{Statement and the parameter freedom used}
The target~\cite{nnlifentry} asks for some $a>0$, $b<0$, $d>0$, real $V_R<V_F$, and $T>0$, for which
\begin{align}
 p_t+\partial_v[(-v+bN(t-d))p]-a p_{vv}&=N(t)\delta_{V_R},\label{eq:physical}\\
 p(t,V_F)=0,\qquad N(t)&=-a p_v(t,V_F),\qquad
 \int_{-\infty}^{V_F}p(t,v)\dd v=1\label{eq:physical-bc}
\end{align}
has a nonnegative periodic solution with nonconstant $N$ and the stated reset, regularity, flux and moment properties.


The Gaussian-wave reduction of Ikeda, Roux, Salort and Smets~\cite{ikeda} and the survey~\cite{crsurvey}, Section 1.7.2, motivate the question. The linear-instability criterion of C\'aceres, Ca\~nizo and Ramos-Lora~\cite{ccr}, Theorem 1.4 and Remark 1.5, does not establish nonlinear periodic existence. Rieutord and Salort~\cite{rs}, Theorem 1 and Corollary 2, study successive large-delay and large-time limits, with positive physical thresholds. Here we construct a local periodic branch at a finite delay, allowing negative thresholds as the repository target permits.

The argument below fixes $a=1$. It does \emph{not} prove a bifurcation for every prescribed pair of physical thresholds. Instead it uses the existential freedom in the target to choose
\begin{equation}\label{eq:threshold-family}
 V_F=bN_0,\qquad V_R=bN_0-1.
\end{equation}
This keeps a convenient centered stationary profile fixed as $b$ varies. Both thresholds will be negative. The catalogue imposes no positivity condition on them.

\begin{theorem}[Proposed existence result]\label{thm:nnlif}
For a sufficiently large frequency $\omega_0$, the construction below gives a fixed phase $\delta\in(0,2\pi)$ and a local nontrivial branch of parameters $b(\kappa)<0$, $\omega(\kappa)>0$. For sufficiently small nonzero real $\kappa$, setting
\[
 a=1,\quad d=\frac{\delta}{\omega(\kappa)},\quad
 T=\frac{2\pi}{\omega(\kappa)},\quad
 V_F=b(\kappa)N_0,\quad V_R=V_F-1
\]
gives a solution of~\eqref{eq:physical}--\eqref{eq:physical-bc} with all the properties required in the target and with a nonzero first firing-rate harmonic.
\end{theorem}

The proof is given rather than inferred from linear instability. The stationary profile, reset resolvent, response asymptotics, nonlinear reduction and positivity are all needed. This manuscript is submitted as a solution claim awaiting independent review.

\section{The centered stationary reset process}
Use a centered voltage $w<0$, with reset at $-1$ and threshold at $0$ (Figure~\ref{fig:reset}).

\begin{figure}[ht]
\centering
\setlength{\unitlength}{1pt}
\begin{picture}(320,83)
\put(15,30){\vector(1,0){285}}
\put(10,12){\makebox(45,10){$-\infty$}}
\put(160,26){\line(0,1){8}}\put(280,24){\line(0,1){12}}
\put(130,10){\makebox(60,10){reset $w=-1$}}
\put(244,10){\makebox(74,10){threshold $w=0$}}
\put(280,42){\line(0,1){22}}\put(160,64){\line(1,0){120}}
\put(160,64){\vector(0,-1){22}}
\put(184,68){\makebox(72,10){reset flux $N$}}
\put(25,44){\makebox(111,10){density domain $w<0$}}
\end{picture}
\caption{The centered voltage domain. Probability reaching the absorbing firing threshold is reinjected at the reset point. The physical coordinate is $v=w+bN_0$, so $V_R=bN_0-1$ and $V_F=bN_0$.}
\label{fig:reset}
\end{figure}
 Define
\begin{align}
 m_0&=\int_{-1}^0 e^{z^2/2}\int_{-\infty}^{z}e^{-y^2/2}\dd y\dd z,
 &N_0&=m_0^{-1},\label{eq:N0}\\
 P(w)&=N_0e^{-w^2/2}\int_{\max(w,-1)}^0e^{z^2/2}\dd z.
 \label{eq:steadyP}
\end{align}
Then $P>0$ for $w<0$, $P(0)=0$, $\int P=1$, $-P'(0)=N_0$, and
\[
 [P']_{-1}:=P'(-1+)-P'(-1-)=-N_0=P'(0).
\]
On each side of $-1$, $P''+(wP)'=0$.

Let
\[
 \HH=L^2(( -\infty,0),e^{w^2/2}\dd w),\qquad
 \HH_0=\left\{q\in\HH:\int_{-\infty}^0q\dd w=0\right\}.
\]
The integral defining $\HH_0$ is bounded on $\HH$ by Cauchy--Schwarz. Complexify these spaces when taking Fourier transforms or resolvents.

Define the reset operator $A$ by
\begin{equation}\label{eq:Adef}
 Aq=q''+(wq)' \quad\hbox{on }(-\infty,-1)\cup(-1,0),
 \qquad Cq=-q'(0).
\end{equation}
Its domain consists of the functions for which $r(w)=e^{w^2/4}q(w)$ is continuous at $-1$, belongs to $H^2$ on both component intervals, satisfies $w^2r\in L^2$, and obeys
\begin{equation}\label{eq:Adomain}
 q(0)=0,\qquad [q']_{-1}=q'(0).
\end{equation}
The corresponding transformed jump is $[r']_{-1}=e^{1/4}r'(0)$. Distributionally the equation $q_t=Aq$ is the absorbing Ornstein--Uhlenbeck equation with reset source $Cq\,\delta_{-1}$: the distributional second derivative contributes $[q']_{-1}\delta_{-1}=-Cq\,\delta_{-1}$, which the source cancels.

Integration on the two component intervals gives
\begin{equation}\label{eq:massA}
 \int Aq\dd w=q'(0)-[q']_{-1}=0.
\end{equation}
Thus $A$ restricts to an operator $A_0$ on $\HH_0$, and $AP=0$.

\section{The reset resolvent and its imaginary-axis bounds}
Let $D$ denote the absorbing Ornstein--Uhlenbeck operator without reset, with Dirichlet boundary condition at zero and without a derivative jump at $-1$. The unitary map $q\mapsto e^{w^2/4}q$ transforms it into
\[
 L=\partial_{ww}+\frac12-\frac{w^2}{4}
\]
on $(-\infty,0)$ with Dirichlet boundary condition. It is self-adjoint, with compact resolvent and eigenvalues $-1,-3,-5,\ldots$.

Write $R_D(\lambda)=(\lambda-D)^{-1}$ and
\[
 g_\lambda=R_D(\lambda)\delta_{-1},\qquad
 a(\lambda)=Cg_\lambda.
\]
The point source is interpreted via the Green function. It belongs to $\HH$ after application of the resolvent and has derivative jump $-1$.

\begin{lemma}[Rank-one reset formula]\label{lem:reset-resolvent}
For $\Rea\lambda\geq0$, $\lambda\ne0$,
\begin{equation}\label{eq:reset-resolvent}
 R_A(\lambda)h=R_D(\lambda)h+
 g_\lambda\frac{CR_D(\lambda)h}{1-a(\lambda)}.
\end{equation}
The operator $A$ is densely defined, closed and has compact resolvent. Its zero eigenvalue is algebraically simple, with eigenfunction $P$. On $\HH_0$ the whole imaginary axis belongs to the resolvent set. There is a constant $c_0$ such that
\begin{equation}\label{eq:imag-bounds}
 \|R_{A_0}(i\xi)\|_{\HH_0\to\HH_0}
 \leq\frac{c_0}{1+|\xi|},\qquad
 \|A_0R_{A_0}(i\xi)\|_{\HH_0\to\HH_0}\leq c_0.
\end{equation}
\end{lemma}
\begin{proof}
If $(\lambda-A)q=h$, then its distributional absorbing equation is
\[
 (\lambda-D)q=h+Cq\,\delta_{-1}.
\]
Consequently $q=R_D(\lambda)h+Cq\,g_\lambda$. Taking $C$ gives
\[
 (1-a(\lambda))Cq=CR_D(\lambda)h,
\]
and hence~\eqref{eq:reset-resolvent} whenever its denominator is nonzero. Conversely that formula has precisely the jump~\eqref{eq:Adomain}, so it solves the original operator equation.

To understand the denominator, let $\tau_0$ be the first hitting time of zero for
\[
 dW_t=-W_t\dd t+\sqrt2\dd B_t,\qquad W_0=-1.
\]
The absorbing flux identity gives $a(\lambda)=\mathbb E_{-1}e^{-\lambda\tau_0}$. One can also verify this identity by solving the adjoint boundary-value equation below. The hitting-time density from $-a$, $a>0$, is
\begin{equation}\label{eq:hittingdensity}
 \rho_a(t)=\frac{2a e^{2t}}{\sqrt{2\pi}(e^{2t}-1)^{3/2}}
 \exp\left[-\frac{a^2}{2(e^{2t}-1)}\right],\qquad t>0.
\end{equation}
Indeed, $W_t=e^{-t}(-a+\widetilde B_{e^{2t}-1})$ for a standard Brownian motion $\widetilde B$. Changing variables in its first-passage density gives~\eqref{eq:hittingdensity}.

The density is positive on every positive time interval. Therefore $|a(i\xi)|<1$ for $\xi\ne0$: equality in the triangle inequality would require $e^{-i\xi t}$ to be constant on the support of this density. For $\Rea\lambda>0$, one similarly has $|a(\lambda)|<1$. At zero,
\[
 a(0)=1,\qquad a'(0)=-\mathbb E_{-1}\tau_0=-m_0<0.
\]
The equality of the mean hitting time and~\eqref{eq:N0} follows either by integrating~\eqref{eq:hittingdensity} or by solving $m''-wm'=-1$, $m(0)=0$, at $w=-1$.

Here are the high-frequency estimates required in the formula. For $|\xi|\geq1$,
\begin{align}
 \|R_D(i\xi)\|&\leq |\xi|^{-1},\label{eq:absorbing1}\\
 \|R_D(i\xi)\delta_{-1}\|_\HH&\leq c|\xi|^{-3/4},\label{eq:absorbing2}\\
 \|CR_D(i\xi)\|_{\HH\to\C}&\leq c|\xi|^{-1/4}.\label{eq:absorbing3}
\end{align}
For detail, the full-line kernel of $e^{tL}$ is
\[
 k_t(v,w)=\frac{e^{t/2}}{\sqrt{4\pi\sinh t}}
 \exp\left[-\frac{(v^2+w^2)\cosh t-2vw}{4\sinh t}\right].
\]
The Dirichlet kernel on the half-line is $k_t(v,w)-k_t(v,-w)$. At an interior diagonal point it is bounded by $ct^{-1/2}$ for $0<t\leq1$. Its mixed boundary derivative at $(0,0)$ is bounded by $ct^{-3/2}$. If $\nu_j=2j+1$ and $\varphi_j$ are normalized Dirichlet eigenfunctions of $-L$, these bounds imply
\[
 \sum_{\nu_j\leq M}|\varphi_j(-1)|^2\leq cM^{1/2},\qquad
 \sum_{\nu_j\leq M}|\varphi_j'(0)|^2\leq cM^{3/2},\quad M\geq1.
\]
For example, apply the respective heat-kernel bound at $t=1/M$ and use $e^{-\nu_j/M}\geq e^{-1}$ in the indicated sum. Integrating the resulting spectral measures against $(\xi^2+\nu^2)^{-1}$, or splitting them into dyadic intervals above and below $|\xi|$, gives squared-norm bounds $c|\xi|^{-3/2}$ and $c|\xi|^{-1/2}$. The harmless factor $e^{1/4}$ from the transformed point source does not change their powers. This proves~\eqref{eq:absorbing2}--\eqref{eq:absorbing3};~\eqref{eq:absorbing1} is immediate from the spectrum.

The density~\eqref{eq:hittingdensity} is integrable, so $a(i\xi)\to0$. Formula~\eqref{eq:reset-resolvent} and~\eqref{eq:absorbing1}--\eqref{eq:absorbing3} now give $\|R_A(i\xi)\|\leq c/|\xi|$ for large $|\xi|$.

For clarity, the domain and closedness assertions can also be checked from the same construction, rather than presupposed. The Dirichlet boundary estimate on nested intervals next to zero gives
\[
 \|r\|_{H^2(-1/2,0)}\leq c\bigl(\|Lr\|_{L^2(-3/4,0)}
                                  +\|r\|_{L^2(-3/4,0)}\bigr).
\]
It follows by the usual cutoff and one-dimensional interpolation estimate, since the potential is bounded there and $r(0)=0$. In particular $|Cq|\leq c(\|Aq\|_\HH+\|q\|_\HH)$. For any fixed positive $\lambda$, the identity
\[
 q=R_D(\lambda)(\lambda q-Aq)+Cq\,g_\lambda
\]
then controls the full piecewise harmonic-oscillator $H^2$ norm by the graph norm. The converse bound follows directly from the definition. Completeness of this piecewise domain, including its continuous trace conditions, proves closedness. Smooth compactly supported functions avoiding $-1$ and $0$ show density. The resolvent formula is a compact absorbing resolvent plus a rank-one operator, hence is compact.

Near zero, the formula has an at-most-simple pole because $1-a(\lambda)=m_0\lambda+O(\lambda^2)$. The kernel is one-dimensional, as is also immediate from that formula at $\lambda=0$. It contains the normalized $P$. From~\eqref{eq:massA},
\[
 \int R_A(\lambda)h\dd w=\lambda^{-1}\int h\dd w.
\]
The residue is therefore $h\mapsto P\int h$. It vanishes on $\HH_0$, so the apparent pole is removable on that space and $-A_0$ is invertible. On compact subsets of the rest of the imaginary axis, the resolvent is continuous and bounded. This proves the first bound in~\eqref{eq:imag-bounds}. The identity $A_0R_{A_0}(i\xi)=i\xi R_{A_0}(i\xi)-I$ proves the second.
\end{proof}

\section{Boundary traces and periodic function spaces}
Equip $D(A)$ with its graph norm. The graph-norm equivalence just proved and the one-dimensional derivative-trace interpolation inequality give
\begin{equation}\label{eq:trace-interpolation}
 |Cq|\leq c\|q\|_\HH^{1/4}\|q\|_{D(A)}^{3/4}.
\end{equation}
To see the exponents directly, apply the local inequality
\[
 |r'(0)|\leq c\|r\|_{L^2(-1/2,0)}^{1/4}
                  \|r\|_{H^2(-1/2,0)}^{3/4}
             +c\|r\|_{L^2(-1/2,0)}
\]
and absorb the last term into the graph norm. The weight is bounded above and below on this interval. Also
\begin{equation}\label{eq:derivative-interpolation}
 \|q_w\|_\HH\leq
 c\|q\|_\HH^{1/2}\|q\|_{D(A)}^{1/2}.
\end{equation}
Indeed, in transformed variables $q_w=e^{-w^2/4}(r'-wr/2)$. Piecewise one-dimensional interpolation controls $r'$, and
$\|wr\|_2^2\leq\|r\|_2\|w^2r\|_2$ controls the second term. Continuity of $r$ at the reset makes its piecewise first derivative its global weak first derivative.

Let $\T=\R/(2\pi\mathbb Z)$, and initially use the real Hilbert spaces
\begin{equation}\label{eq:XY}
 X=H^1(\T;\HH_0),\qquad
 Y=H^2(\T;\HH_0)\cap H^1(\T;D(A_0)).
\end{equation}
For every $\omega>0$, Lemma~\ref{lem:reset-resolvent}, applied to Fourier coefficients, shows that
\begin{equation}\label{eq:periodiciso}
 \omega\partial_\theta-A_0:Y\longrightarrow X
 \quad\hbox{is an isomorphism}.
\end{equation}
More generally it is an isomorphism
\[
 Y_s:=H^{s+1}(\T;\HH_0)\cap H^s(\T;D(A_0))
 \longrightarrow H^s(\T;\HH_0)
\]
for every $s\geq0$. The $n$th inverse multiplier is $R_{A_0}(in\omega)$; the two bounds in~\eqref{eq:imag-bounds} give exactly one time derivative in $\HH_0$ and one graph norm without a time loss.

Interpolation~\eqref{eq:trace-interpolation} on each Fourier coefficient and H\"older's inequality in the frequency sum give
\begin{equation}\label{eq:time-trace}
 Cq\in H^{5/4}(\T),\qquad
 \|Cq\|_{H^{5/4}}\leq c\|q\|_Y.
\end{equation}
Explicitly, the weight $\langle n\rangle^{5/2}$ on $|Cq_n|^2$ is the product of the $\langle n\rangle^4\|q_n\|_\HH^2$ weight to the power $1/4$ and the $\langle n\rangle^2\|q_n\|_{D(A)}^2$ weight to the power $3/4$. Thus $q\mapsto Cq$ is compact from $Y$ to $H^1(\T)$. Likewise~\eqref{eq:derivative-interpolation} gives $q_w\in H^{3/2}(\T;\HH)$; the weaker inclusion $H^1(\T;\HH)$ already suffices for the nonlinear map below.

\section{An explicit firing-response function}
Put
\begin{equation}\label{eq:response}
 f=-P'\in\HH_0,
 \qquad K(\lambda)=C(\lambda-A_0)^{-1}f.
\end{equation}
For the imaginary-axis values, the resolvent has just been justified. Let
\[
 h_\lambda(w)=\mathbb E_w e^{-\lambda\tau_0},\qquad w<0,
\]
where the absorbing OU process is now started at $w$. It solves
\begin{equation}\label{eq:backward}
 h_\lambda''-w h_\lambda'=\lambda h_\lambda,
 \qquad h_\lambda(0)=1.
\end{equation}
An integration by parts in the absorbing resolvent equation gives
\[
 CR_D(\lambda)g=\int_{-\infty}^0h_\lambda(w)g(w)\dd w,
 \qquad a(\lambda)=h_\lambda(-1).
\]
These identities hold first for $\Rea\lambda>0$ and extend to nonzero imaginary $\lambda$.

\begin{lemma}[Response formula and phase]\label{lem:response}
For $\lambda=i\xi\ne0$,
\begin{equation}\label{eq:Kexact}
 K(\lambda)=
 \frac{N_0\{h_\lambda'(0)-h_\lambda'(-1)\}}
 {(\lambda+1)\{1-h_\lambda(-1)\}}.
\end{equation}
As $\xi\to+\infty$,
\begin{align}
 K(i\xi)&=N_0(i\xi)^{-1/2}
 \left(1-\frac{5}{4i\xi}+O(\xi^{-2})\right),\label{eq:Kasympt}\\
 \frac{\mathrm d}{\mathrm d\xi}\arg K(i\xi)
 &=-\frac{5}{4\xi^2}+O(\xi^{-3}).\label{eq:phase-derivative}
\end{align}
The argument is chosen continuously for sufficiently large positive $\xi$. Finally, $K(0)>0$.
\end{lemma}
\begin{proof}
The reset formula yields
\[
 K(\lambda)=\frac{\int h_\lambda(-P')\dd w}{1-h_\lambda(-1)}
 =\frac{\int h_\lambda'P\dd w}{1-h_\lambda(-1)}.
\]
For any admissible test function $\varphi$, stationarity of $P$ gives
\[
 \int P(\varphi''-w\varphi')\dd w
 +N_0\{\varphi(-1)-\varphi(0)\}=0.
\]
Differentiating~\eqref{eq:backward} shows
\[
 (\partial_{ww}-w\partial_w)h_\lambda'=(\lambda+1)h_\lambda'.
\]
Taking $\varphi=h_\lambda'$ proves~\eqref{eq:Kexact}. Gaussian decay of $P$ justifies the lower-end integrations; the hitting-time representation also bounds the functions involved there.

For $\Rea\lambda>0$, an integral solution of~\eqref{eq:backward} is
\[
 h_\lambda(w)=
 \frac{\int_0^\infty s^{\lambda-1}e^{-s^2/2+ws}\dd s}
      {\int_0^\infty s^{\lambda-1}e^{-s^2/2}\dd s}.
\]
Integration by parts in $s$ verifies the differential equation, while the boundary value and decay at $-\infty$ select the hitting-time solution. At zero this gives, by continuation,
\begin{equation}\label{eq:gamma-ratio}
 h_\lambda'(0)=\sqrt2\,
 \frac{\Gamma((\lambda+1)/2)}{\Gamma(\lambda/2)}
 =\lambda^{1/2}\left(1-\frac1{4\lambda}+O(\lambda^{-2})\right).
\end{equation}
The last expansion is the gamma-ratio asymptotic~\cite{dlmf}, valid uniformly in a sector containing the positive imaginary axis.

For fixed starting distance $a>0$, the density~\eqref{eq:hittingdensity}, all its time derivatives, and its derivatives in $a$ are flat at $t=0$ and decay exponentially as $t\to\infty$. Multiplication by any fixed power of $t$ preserves integrability of all the required derivatives. Repeated integration by parts in $t$ therefore gives, for every integer $m$,
\[
 h_{i\xi}(-1)=O_m(\xi^{-m}),\qquad
 h_{i\xi}'(-1)=O_m(\xi^{-m}),
\]
with the same conclusion after differentiation in $\xi$. In~\eqref{eq:Kexact}, these remote-reset contributions are thus smaller than every inverse power. Combining~\eqref{eq:gamma-ratio} with
$(\lambda+1)^{-1}=\lambda^{-1}(1-\lambda^{-1}+O(\lambda^{-2}))$
proves~\eqref{eq:Kasympt}. The coefficient $-5/4$ is the sum of $-1/4$ from the gamma ratio and $-1$ from this last factor. The sector-uniform analytic remainder in~\eqref{eq:gamma-ratio} may be differentiated by Cauchy estimates on discs of radius proportional to $|\lambda|$ inside the sector. Combined with the differentiated hitting-time estimates, this gives
\[
 \arg K(i\xi)=-\frac\pi4+\frac5{4\xi}+O(\xi^{-2}),
\]
and hence~\eqref{eq:phase-derivative}.

To check the sign at zero, add a constant centered input $I$ to the drift while keeping the centered thresholds $-1,0$ fixed. The mean inter-firing time and stationary firing rate are
\[
 m(I)=\int_{-1}^0\int_0^\infty e^{(z-I)s-s^2/2}\dd s\dd z,
 \qquad N(I)=m(I)^{-1}.
\]
Clearly $m'(I)<0$ and therefore $N'(0)>0$. The normalized stationary density $P_I$ is differentiable in the graph norm near zero, as is seen from its explicit Gaussian formula. Differentiating $AP_I-I(P_I)'=0$ at $I=0$ and imposing zero derivative of the mass gives
\[
 (-A_0)\dot P_0=-P'=f,
 \qquad K(0)=C\dot P_0=N'(0)>0.
\]
\end{proof}

\section{Choosing a simple pair of critical harmonics}
Choose $\omega_0$ sufficiently large that $K(i\omega_0)\ne0$ and the derivative in~\eqref{eq:phase-derivative} is nonzero. Also arrange
\begin{equation}\label{eq:harmonic-separation}
 |K(in\omega_0)|<|K(i\omega_0)|\qquad(n\geq2).
\end{equation}
This last choice is uniform over the infinitely many harmonics. Indeed, for sufficiently large $\Omega$ and all $\xi\geq\Omega$,~\eqref{eq:Kasympt} gives
\[
 0.9N_0\xi^{-1/2}\leq|K(i\xi)|\leq1.1N_0\xi^{-1/2}.
\]
The ratio in~\eqref{eq:harmonic-separation} is at most
$(1.1/0.9)n^{-1/2}<1$ for every $n\geq2$.

Define
\begin{equation}\label{eq:criticalparams}
 b_0=-\frac1{|K(i\omega_0)|}<0,
 \qquad b_0e^{-i\delta}K(i\omega_0)=1.
\end{equation}
The phase can be selected in $(0,2\pi)$; asymptotically it is $3\pi/4+O(\omega_0^{-1})$. From this point onward, \emph{$\delta$ is fixed}. The physical delay will be $d=\delta/\omega$, not an independently translated variable in the Banach-space map.

\section{The full nonlinear periodic equation}
Write the centered density as $P(w)+q(\theta,w)$, with $\theta=\omega t$, and put $n(\theta)=Cq(\theta)$. Under~\eqref{eq:threshold-family}, the physical drift in centered coordinates is
\[
 -w+b\{N(t-d)-N_0\}=-w+b n(\theta-\delta).
\]
The full periodic equation is therefore
\begin{equation}\label{eq:nonlinearF}
 F(q,b,\omega):=\omega q_\theta-A_0q
       +b(Cq)(\theta-\delta)(P'+q_w)=0.
\end{equation}
It includes the original nonlinear feedback, not only its linearization.

The map $F:Y\times\R\times(0,\infty)\to X$ is smooth (in fact real analytic). To verify the only nontrivial mapping assertion, $Cq\in H^{5/4}(\T)\subset H^1(\T)$ and $q_w\in H^1(\T;\HH)$. The one-dimensional vector-valued $H^1$ product inequality bounds their product in $H^1(\T;\HH)$. Both $\int P'=0$ and $\int q_w=0$, so the forcing has zero mass. Translation by the fixed $\delta$ is a bounded linear operator on every time Sobolev space. No differentiability of a variable-delay translation is being assumed.

At $(0,b_0,\omega_0)$ the derivative in $q$ is
\[
 \mathcal Lq=\omega_0q_\theta-A_0q
             -b_0f(Cq)(\theta-\delta).
\]
It is a compact perturbation of~\eqref{eq:periodiciso}, by the compact trace map following~\eqref{eq:time-trace}; hence it is Fredholm of index zero. The $n$th Fourier equation has characteristic factor
\begin{equation}\label{eq:charfactor}
 D_n=1-b_0e^{-in\delta}K(in\omega_0).
\end{equation}
Only $n=1,-1$ vanish. For $n=0$ this follows from $b_0<0<K(0)$, and for $|n|\geq2$ from~\eqref{eq:harmonic-separation} and complex conjugation.

Let
\[
 R=(i\omega_0-A_0)^{-1},\qquad
 e=\frac{Rf}{K(i\omega_0)},\qquad Ce=1,
 \qquad \ell=CR.
\]
The real kernel is spanned by the real and imaginary parts of $e^{i\theta}e$. Moreover, $\ell$ annihilates the first Fourier coefficient of every $\mathcal Lq$ and $\ell(f)=K(i\omega_0)\ne0$. The range is therefore exactly
\[
 Z=\{h\in X:\ell(h_1)=0\},
\]
where a complex equation denotes two real conditions. Indeed, this subspace has real codimension two, and the Fredholm index and the two-dimensional kernel give the same codimension for the range. On the complementary domain
\[
 Y_c=\{q\in Y:Cq_1=0\},
\]
$\mathcal L:Y_c\to Z$ is an isomorphism.

Choose the bounded projection $\Pi:X\to Z$ that subtracts
\[
 e^{i\theta}f\,\frac{\ell(h_1)}{K(i\omega_0)}
 +e^{-i\theta}f\,\overline{\frac{\ell(h_1)}{K(i\omega_0)}}
\]
from a real $h$. Seek
\begin{equation}\label{eq:branchansatz}
 q=\kappa\{e^{i\theta}e+e^{-i\theta}\bar e\}+y,
 \qquad y\in Y_c,\quad\kappa\in\R.
\end{equation}
The implicit function theorem applied to $\Pi F=0$ gives a smooth
$y=y(\kappa,b,\omega)$ near $(0,b_0,\omega_0)$, with
\[
 y(0,b,\omega)=0,\qquad y_\kappa(0,b_0,\omega_0)=0.
\]
The remaining equation is the single complex equation
\[
 g(\kappa,b,\omega):=\ell(F(q,b,\omega)_1)=0.
\]
Since $g(0,b,\omega)=0$, division by $\kappa$ extends smoothly to a function $G$ at $\kappa=0$. At the base point $G=0$.

The real two-parameter Jacobian of $G$ is nonsingular. Its two complex column vectors are
\begin{align}
 \partial_bG&=\ell(-e^{-i\delta}f)=-\frac1{b_0},\label{eq:jacb}\\
 \partial_\omega G&=i\,\frac{CR^2f}{K(i\omega_0)}
          =-i\,\frac{K'(i\omega_0)}{K(i\omega_0)}.\label{eq:jacomega}
\end{align}
Terms involving derivatives of $y$ disappear because $\ell$ annihilates the range of $\mathcal L$. In~\eqref{eq:jacomega}, the prime denotes differentiation in the complex resolvent variable. The first column is real and nonzero, while
\[
 \Ima\partial_\omega G
 =-\Rea\frac{K'(i\omega_0)}{K(i\omega_0)}
 =-\left.\frac{\mathrm d}{\mathrm d\omega}\arg K(i\omega)
   \right|_{\omega=\omega_0}\ne0.
\]
The real implicit function theorem thus supplies $b=b(\kappa)$ and $\omega=\omega(\kappa)$ such that $F=0$ for small $\kappa$. These parameters stay negative and positive, respectively. By the defining complement in~\eqref{eq:branchansatz},
\begin{equation}\label{eq:nonconstant}
 (Cq)_1=\kappa.
\end{equation}
Every nonzero branch parameter gives a nonconstant firing rate.

\section{Positivity is a separate step, not an interior-cone assumption}
Smallness in the Gaussian Hilbert norm alone does not guarantee $P+q\geq0$ in the far tail. We prove positivity directly.

For a sufficiently small branch parameter,~\eqref{eq:time-trace} and the time Sobolev embedding give
\begin{equation}\label{eq:smallinput}
 N(\theta):=N_0+Cq(\theta)\geq N_0/2,
 \qquad |I(\theta)|:=|b(Cq)(\theta-\delta)|\leq1/2.
\end{equation}
Set $r=e^{w^2/4}(P+q)$ and $z=r_-:=\max(-r,0)$. The transformed distributional equation is
\begin{equation}\label{eq:rPDE}
 \omega r_\theta=r_{ww}+(1/2-w^2/4)r
 -I(\theta)(r_w-wr/2)+e^{1/4}N(\theta)\delta_{-1}.
\end{equation}
The available graph regularity makes $r$ globally $H^1$ in $w$, with the prescribed derivative jump, and makes the negative-part test legitimate. Its trace is zero at $w=0$. Multiplying~\eqref{eq:rPDE} by $-z$ and integrating gives
\begin{equation}\label{eq:negativeenergy}
 \frac\omega2\frac{\mathrm d}{\mathrm d\theta}\|z\|_2^2
 \leq-\|z_w\|_2^2+\frac12\|z\|_2^2
       -\frac14\|wz\|_2^2+
       \frac{I(\theta)}2\int w z^2\dd w.
\end{equation}
The reset contribution is $-e^{1/4}N(\theta)z(-1)\leq0$; the constant-in-$w$ transport contribution $-I r_w$ integrates to zero on the negative set.

The Dirichlet half-line oscillator inequality is
\begin{equation}\label{eq:oscillatorineq}
 \|z_w\|_2^2+\tfrac14\|wz\|_2^2\geq\tfrac32\|z\|_2^2.
\end{equation}
It follows from the first odd oscillator eigenfunction $(-w)e^{-w^2/4}$, or directly by its ground-state factorization. Young's inequality gives
\[
 \left|\frac I2\int wz^2\dd w\right|
 \leq\frac1{16}\|wz\|_2^2+I^2\|z\|_2^2.
\]
Subtracting only three quarters of the oscillator form in~\eqref{eq:negativeenergy} leaves a further nonpositive derivative term. Using~\eqref{eq:oscillatorineq} and~\eqref{eq:smallinput},
\[
 \frac\omega2\frac{\mathrm d}{\mathrm d\theta}\|z\|_2^2
 \leq\left(-\frac58+I^2\right)\|z\|_2^2
 \leq-\frac38\|z\|_2^2.
\]
Integrating over a full period forces $z=0$. Thus $P+q\geq0$ for every time.

\section{Regularity, flux, moment and the original variables}
The branch initially lies in $Y=Y_1$. It can be bootstrapped without assuming a semigroup theorem for the nonlinear delay equation. If $q\in Y_s$ with $s\geq1$, the same Fourier interpolation used above gives
\[
 Cq\in H^{s+1/4}(\T),\qquad
 q_w\in H^{s+1/2}(\T;\HH).
\]
The one-dimensional Sobolev multiplication rule implies that the right side of
\[
 (\omega\partial_\theta-A_0)q
 =-b(Cq)(\theta-\delta)(P'+q_w)
\]
belongs to $H^{s+1/4}(\T;\HH_0)$. The periodic inverse then gives $q\in Y_{s+1/4}$. Iteration gives arbitrarily high time regularity in the graph domain. On each open voltage interval away from the reset, the differential equation then gives arbitrarily high local spatial regularity by ordinary elliptic bootstrapping, applied also to time derivatives. Thus the density is classical away from the reset, continuous at the reset, and has exactly its prescribed derivative jump.

The mass is one because $q(\theta)\in\HH_0$. Uniform boundedness of $P+q$ in $\HH$ implies
\[
 \sup_\theta\int_{-\infty}^0(1+w^2)|P(w)+q(\theta,w)|\dd w<\infty
\]
by Cauchy--Schwarz with the Gaussian reciprocal weight. For the flux, the transformed graph domain has $r$ and $r_w$ bounded on the lower half-line; consequently
\[
 (-w+I)(P+q)-(P+q)_w
 =e^{-w^2/4}\{(-w/2+I)r-r_w\}\longrightarrow0
 \quad(w\to-\infty).
\]
This verifies the lower-boundary flux condition.

Finally, for a fixed small nonzero branch parameter define
\begin{equation}\label{eq:reconstruction}
 p(t,v)=P(v-bN_0)+q(\omega t,v-bN_0),
 \qquad N(t)=N_0+Cq(\omega t).
\end{equation}
With~\eqref{eq:threshold-family}, $d=\delta/\omega$ and $T=2\pi/\omega$, the drift identity preceding~\eqref{eq:nonlinearF} shows that~\eqref{eq:reconstruction} solves the original equations, not a forced or reduced model. The fixed finite voltage translation preserves the second-moment bound. The flux identity at the firing boundary, the reset jump and periodicity are already built into the construction. Nonconstancy follows from~\eqref{eq:nonconstant}. This completes the proposed proof of Theorem~\ref{thm:nnlif}.

\begin{remark}[What this construction does not establish]
It does not assert stability of the periodic branch, uniqueness, existence for a prescribed biological parameter set, or a large-delay asymptotic description. None of these is needed for the existential catalogue statement. The small amplitude is selected after the large but finite $\omega_0$ and $b_0$, so the requirements $|b n|\leq1/2$ and $N_0+n>0$ are compatible.
\end{remark}


\begin{thebibliography}{99}
\bibitem{nnlifentry} M. Colbrook, \emph{AIM 353: A genuinely periodic firing pattern in the delayed noisy integrate-and-fire PDE}, repository revision \texttt{fa98b7525fa3}, consulted 2 October 2026. \href{https://github.com/MColbrook/AIM/blob/fa98b7525fa3f78317536a8825f9cfa0ae1c369c/problems/353-delayed-nnlif-periodic.md}{Pinned problem statement}.

\bibitem{ikeda} K. Ikeda, P. Roux, D. Salort and D. Smets, \emph{Theoretical study of the emergence of periodic solutions for the inhibitory NNLIF neuron model with synaptic delay}, Mathematical Neuroscience and Applications 2 (2022), article 7256. \href{https://mna.episciences.org/10212}{Publisher version}. \href{https://doi.org/10.46298/mna.7256}{doi:10.46298/mna.7256}. The periodicity result concerns the Gaussian-wave reduction.

\bibitem{ccr} M. J. C\'aceres, J. A. Ca\~nizo and A. Ramos-Lora, \emph{On the asymptotic behavior of the NNLIF neuron model for general connectivity strength}, Communications in Mathematical Physics 406 (2025), article 115. \href{https://arxiv.org/html/2401.13534v1}{arXiv:2401.13534v1}. \href{https://doi.org/10.1007/s00220-025-05287-5}{doi:10.1007/s00220-025-05287-5}. Linear stability and instability are distinct from nonlinear periodic existence.

\bibitem{rs} C. Rieutord and D. Salort, \emph{Asymptotic dynamics of inhibitory networks for the NNLIF Model in the large-delay limit}, \href{https://arxiv.org/html/2606.17611v2}{arXiv:2606.17611v2}, 2026. The finite-window large-delay result is not a periodic orbit at fixed finite delay.

\bibitem{crsurvey} J. A. Carrillo and P. Roux, \emph{Nonlinear partial differential equations in neuroscience: from modelling to mathematical theory}, Mathematical Models and Methods in Applied Sciences 35 (2025), 403--584. \href{https://arxiv.org/html/2501.06015v1}{arXiv:2501.06015v1}, \href{https://doi.org/10.1142/S0218202525400044}{doi:10.1142/S0218202525400044}. Section 1.7.2.

\bibitem{dlmf} NIST Digital Library of Mathematical Functions, \href{https://dlmf.nist.gov/5.11.iii}{Section 5.11(iii): Ratios}, equations (5.11.13)--(5.11.15), consulted 2 October 2026.
\end{thebibliography}
\end{document}
